\documentclass[12pt]{article}
\usepackage[english]{babel}
\usepackage[utf8]{inputenc}
\usepackage[T2A]{fontenc}
\usepackage{amssymb}
\usepackage{hyperref, fancyhdr, lastpage, fancyvrb, tcolorbox, titlesec}
\usepackage{array, tabularx, colortbl}
\usepackage{tikz}
\usepackage{venndiagram}
\usepackage{amsthm, bm}
\usepackage{relsize}
\usepackage{amsmath,physics}
\usepackage{mathtools}
\usepackage{subcaption}
\usepackage{scalerel}
\usepackage{theoremref}
\usepackage{circuitikz}
\usepackage{geometry}
\usepackage{stmaryrd}
\usepackage{forest}
\usepackage{cancel}
\usepackage{faktor}
\usepackage{gensymb}
\usetikzlibrary{automata, arrows, positioning, shapes}
\usetikzlibrary{decorations.pathmorphing}
\useforestlibrary{linguistics}

\ExplSyntaxOn
\NewDocumentCommand{\opair}{m}
 {
  \langle\mspace{2mu}
  \clist_set:Nn \l_tmpa_clist { #1 }
  \clist_use:Nn \l_tmpa_clist {,\mspace{3mu plus 1mu minus 1mu}\allowbreak}
  \mspace{2mu}\rangle
}
\ExplSyntaxOff

\hypersetup{
    colorlinks=true,
    linktoc=all,
    linkcolor=blue
}

% \setlength\parindent{0pt}

\newcommand{\N}{\mathbb{N}}
\newcommand{\Z}{\mathbb{Z}}
\newcommand{\Q}{\mathbb{Q}}
\newcommand{\R}{\mathbb{R}}
\newcommand{\E}{\mathbb{E}}

\newcommand{\vars}{\operatorname{Vars}}
\newcommand{\free}{\operatorname{Free}}

\newcommand{\logand}{\; \& \;}

\newcommand{\calA}{\mathcal{A}}
\newcommand{\calS}{\mathcal{S}}
\newcommand{\calD}{\mathcal{D}}
\newcommand{\calL}{\mathcal{L}}
\newcommand{\calZ}{\mathcal{Z}}
\newcommand{\calQ}{\mathcal{Q}}
\newcommand{\calR}{\mathcal{R}}
\newcommand{\calN}{\mathcal{N}}
\newcommand{\calM}{\mathcal{M}}
\newcommand{\calF}{\mathcal{F}}
\newcommand{\calP}{\mathcal{P}}
\newcommand{\calC}{\mathcal{C}}
\newcommand{\calG}{\mathcal{G}}
\newcommand{\calB}{\mathcal{B}}

\newcommand{\continuum}{\mathfrak{c}}

\newcommand{\cof}{\operatorname{cof}}
\newcommand{\dom}{\operatorname{Dom}}
\newcommand{\rng}{\operatorname{Rng}}

\newcommand{\MA}{\mathsf{MA}}

\newcommand{\fa}{\kern.05em\ensuremath\forall\kern-.22em\rotatebox{69}{\rule{.73em}{.4pt}}}
\newcommand{\ex}{\kern.05em\ensuremath\exists\kern-.03em\rotatebox{90}{\rule{.68em}{.35pt}}}

\newcommand{\dequiv}{\stackrel{\text{деф.}}{\longleftrightarrow}}

\newcommand{\db}[1]{\llbracket #1 \rrbracket}

\DeclareMathOperator*{\bigand}{\scalerel*{\&}{\prod}}
\DeclareMathOperator*{\bigor}{\scalerel*{\lor}{\sum}}

\newtheorem*{theorem}{Theorem}
\newtheorem*{claim}{Claim}
% \newtheorem*{definition}{Дефиниция}
% \newtheorem*{claim}{Твърдение}
% \newtheorem*{property}{Свойство}
% \newtheorem*{hint}{Упътване}
% \theoremstyle{definition}
% \newtheorem*{solution}{Решение}
% \newtheorem*{notation}{Нотация}
% \newtheorem{problem}{Задача}

\title{Martin--Solovay Construction}
\author{Todor Dukov}

\begin{document}
\maketitle

\begin{theorem}[Martin--Solovay]
  Martin's Axiom (or $\MA$ for short) implies that $2^\kappa = \mathfrak{c}$ for all infinite cardinals $\kappa < \mathfrak{c}$, and consequently that $\mathfrak{c}$ is regular.
\end{theorem}

\begin{proof}
  Assuming $\MA$, we first show that $2^\kappa \le \mathfrak{c}$ for every infinite cardinal $\kappa < \mathfrak{c}$. 

  Fix an infinite cardinal $\kappa < \mathfrak{c}$.
  Since there exists an almost disjoint family of size $\mathfrak{c}$, we can select an almost disjoint family $\calA = \{ A_\alpha\}_{\alpha < \kappa}$ of size $\kappa$.

  For any $X \subseteq \kappa$, our goal is to construct a set $B_X \subseteq \omega$ such that for all $\alpha < \kappa$,
  \begin{align*}
    \alpha \in X \iff B_X \cap A_\alpha \text{ is infinite.} \tag{$\ast$} \label{coding}
  \end{align*}
  If every subset $X \subseteq \kappa$ is uniquely determined by such a $B_X \subseteq \omega$ via
  \[
    X = \{ \alpha < \kappa \mid B_X \cap A_\alpha \text{ is infinite} \},
  \]
  it yields a surjection from $\mathcal{P}(\omega)$ onto $\mathcal{P}(\kappa)$. Thus, $2^\kappa \le 2^{\aleph_0} = \mathfrak{c}$.

  To construct $B_X$ for a fixed $X \subseteq \kappa$ and almost disjoint family $\calA$, define the partial order $(P_\calA^X, \le_\calA^X)$ as follows:
  \begin{itemize}
    \item Conditions are functions $p$ satisfying:
      \begin{itemize}
        \item $p \subseteq \omega \times \{0, 1\}$,
        \item $\mathrm{dom}(p) \cap A_\alpha$ is finite for every $\alpha \in X$,
        \item $p_+ := \{ n \in \mathrm{dom}(p) \mid p(n) = 1 \}$ is finite.
      \end{itemize}
    \item The ordering is reverse inclusion ($p \le_\calA^X q \iff q \subseteq p$).
  \end{itemize}

  Given a family $\calD$ of dense sets, every $\calD$-generic filter $G \subseteq P_\calA^X$ yields a corresponding subset $B_X := \bigcup \{ p_+ \mid p \in G \}$ of $\omega$.
  To ensure $B_X$ satisfies \eqref{coding}, we construct $\mathcal{D}$ from the following dense sets for each $\alpha < \kappa$:
  \begin{itemize}
    \item For $\alpha \in X$ and $n \in \omega$, let $D'_{\alpha, n} := \{ p \in P_\calA^X \mid p_+ \cap A_\alpha \text{ is of size at least } n \}$. Meeting these sets forces $B_X \cap A_\alpha$ to be infinite.
    \item For $\alpha \notin X$, let $D''_\alpha := \{ p \in P_\calA^X \mid A_\alpha \subseteq \mathrm{dom}(p) \}$. Meeting these sets forces $B_X \cap A_\alpha$ to be finite.
  \end{itemize}
  Since $|\calD| = \kappa < \continuum$, $\MA$ will guarantee the existence of such a $\calD$-generic filter $G$, provided we show that our poset satisfies the countable chain condition and that $\calD$ is indeed a family of dense sets.

  To verify the ccc, note that if $p_+ = q_+$, then $p$ and $q$ are compatible (since $p \cup q$ forms a valid condition).
  Thus, any antichain in our poset can contain at most one condition for each choice of $p_+$.
  Since $p_+$ is a finite subset of $\omega$, there are only countably many such choices, so $(P_\calA^X, \leq_\calA^X)$ is ccc.
  To verify density:
  \begin{itemize}
    \item For $D'_{\alpha, n}$ (where $\alpha \in X$ and $n \in \omega$):
          Given $p \in P_\calA^X$, since $\dom(p) \cap A_\alpha$ is finite, $A_\alpha \setminus \dom(p)$ is infinite.
          Pick distinct $k_1, \dots, k_n \in A_\alpha \setminus \dom(p)$ and extend $p$ to $p'$ by setting $p'(k_1) = \dots = p'(k_n) = 1$.
          Since only finitely many values were added, the extension $p'$ is a valid condition in $P_\calA^X$ belonging to $D'_{\alpha, n}$.
    \item For $D''_\alpha$ (where $\alpha \notin X$):
          Given $p \in P_\calA^X$, extend $p$ to $p'$ by setting $p'(k) = 0$ for all $k \in A_\alpha \setminus \dom(p)$.
          Since $p'_+ = p_+$ and $\dom(p')$ extends $A_\alpha$ by only finitely many elements, $p'$ remains a valid condition in $P_\calA^X$ with $A_\alpha \subseteq \dom(p')$, meaning $p' \in D''_\alpha$.
  \end{itemize}
  
  Let $G \subseteq P_\calA^X$ be the $\calD$-generic filter guaranteed by $\MA$.
  We verify that $B_X$ satisfies \eqref{coding}:
  \begin{itemize}
    \item[$(\Rightarrow)$] 
      Let $\alpha \in X$. 
      For each $n \in \omega$, $G$ intersects $D'_{\alpha, n}$ in some condition $p$, so $p_+ \cap A_\alpha$ has at least $n$ elements.
      Since $p_+ \subseteq B_X$, it follows that $B_X \cap A_\alpha$ is of size at least $n$ for all $n \in \omega$, so $B_X \cap A_\alpha$ is infinite.
    \item[$(\Leftarrow)$] 
      Let $\alpha \in \kappa \setminus X$.
      Since $G$ intersects $D''_\alpha$, pick $p \in G \cap D''_\alpha$.
      Then $A_\alpha \subseteq \dom(p)$, and since all conditions in $G$ are compatible with $p$, no element of $G$ can set $1$ on $A_\alpha \setminus p_+$.
      Thus $B_X \cap A_\alpha \subseteq p_+$, which is finite.
  \end{itemize}

  Finally, $\continuum$ is regular, since
  \[
    \cof(\continuum) = \cof(2^\kappa) > \kappa \text{ for all } \kappa < \continuum.
  \]
\end{proof}

\end{document}
